\chapter{Construction of Brownian Motion}

Brownian motion can be characterized as a centered Gaussian process whose covariance
structure is given by
\[
\E[B(s)B(t)] = \min(s,t), \qquad s,t \in [0,1].
\]
The guiding principle behind the construction below is to \emph{build such a Gaussian
process explicitly} by expanding the indicator functions
$1_{[0,t]}$ in an orthonormal basis of $L^2([0,1])$.

Indeed, if $\{g_{k,n}\}$ is an orthonormal basis of $L^2([0,1])$, then any function
$h \in L^2([0,1])$ admits the expansion
\[
h = \sum_{k,n} \langle h, g_{k,n} \rangle g_{k,n},
\]
and Parseval's identity ensures that inner products are preserved by summation over the
basis. Applying this idea formally to the family $\{1_{[0,t]}\}_{t\in[0,1]}$, one is led
to consider random series of the form
\[
B(t) = \sum_{k,n} Z_{k,n} \int_0^t g_{k,n}(u)\,du,
\]
where $\{Z_{k,n}\}$ are independent standard normal random variables.

The choice of the Haar basis is particularly natural: its elements are localized on
dyadic intervals, and their integrals are continuous, piecewise linear ``tent''
functions. This localization allows for precise control of the uniform convergence of
the series and ultimately guarantees the continuity of the sample paths.

The proof below follows this strategy rigorously. First, we show that the random series
converges uniformly almost surely, producing a continuous process. Then we identify the
finite-dimensional distributions and compute the covariance via Parseval's identity.
The Gaussian structure then automatically yields independent, stationary increments,
revealing the constructed process to be Brownian motion.

\begin{theorem}[Existence (L\'evy--Ciesielski / Haar--Schauder Construction)]
There exists a probability space $(\Omega,\mathcal F,\Prob)$ and a stochastic process
$B=(B_t)_{t\in[0,1]}$ such that

\begin{enumerate}
\item $B_0=0$ almost surely,
\item $t\mapsto B_t(\omega)$ is continuous on $[0,1]$ for $\Prob$-almost every $\omega$,
\item for all $0\le s<t\le 1$, $B_t-B_s\sim\mathcal N(0,t-s)$,
\item increments over disjoint intervals are independent.
\end{enumerate}

Hence $B$ is a standard Brownian motion on $[0,1]$.

Moreover, $B$ admits the series representation
\[
B(t)=\sum_{n=0}^\infty\sum_{k=1}^{2^n} Z_{k,n}\, f_{k,n}(t),
\]
where $\{Z_{k,n}\}$ are independent standard normal random variables and
$f_{k,n}(t)=\int_0^t g_{k,n}(u)\,du$, with $\{g_{k,n}\}$ the Haar orthonormal basis of
$L^2([0,1])$.
\end{theorem}

\begin{proof}

The proof proceeds in four steps:
(i) definition of the deterministic basis,
(ii) almost sure uniform convergence of the random series,
(iii) identification of the Gaussian covariance,
(iv) verification of the Brownian motion properties.

\Step{1: Haar functions and Schauder primitives}

Let $\{g_{k,n}\}_{n\ge0,\,1\le k\le2^n}$ be the Haar orthonormal basis of $L^2([0,1])$,
defined by
\[
g_{k,n}(t)
=
2^{n/2}\big(1_{I_{k,n}^\ell}(t)-1_{I_{k,n}^r}(t)\big),
\]
where $I_{k,n}=[(k-1)2^{-n},k2^{-n})$ and $I_{k,n}^\ell$, $I_{k,n}^r$ are its left and right
halves.

Define the integrated (Schauder) functions
\[
f_{k,n}(t)=\int_0^t g_{k,n}(u)\,du.
\]
Each $f_{k,n}$ is continuous, piecewise linear, and supported on $I_{k,n}$.
A direct computation shows
\[
\|f_{k,n}\|_\infty \le 2^{-(n/2+1)}.
\]

\Step{2: Definition of the random series}

Let $\{Z_{k,n}\}$ be independent $\mathcal N(0,1)$ random variables.
For $N\in\mathbb{N}$ define the partial sums
\[
B_N(t)=\sum_{n=0}^N\sum_{k=1}^{2^n} Z_{k,n} f_{k,n}(t).
\]
Each $B_N$ is a continuous stochastic process on $[0,1]$.

\Step{3: Almost sure uniform convergence}

Define the level-$n$ contribution
\[
\Delta_n(t)=\sum_{k=1}^{2^n} Z_{k,n} f_{k,n}(t).
\]
Since the supports of $f_{k,n}$ are disjoint for fixed $n$,
\[
\sup_{t\in[0,1]}|\Delta_n(t)|
\le \max_{1\le k\le2^n}|Z_{k,n}|\cdot \|f_{k,n}\|_\infty
\le 2^{-(n/2+1)} \max_{k}|Z_{k,n}|.
\]

Using Gaussian tail estimates and the union bound, one shows that
\[
\sum_{n=0}^\infty \Prob\!\left(\max_{k}|Z_{k,n}| > a_n\right)<\infty
\]
for a suitable sequence $a_n\to\infty$.
By the Borel--Cantelli lemma,
\[
\sum_{n=0}^\infty \sup_{t\in[0,1]}|\Delta_n(t)| < \infty
\quad\text{almost surely}.
\]
Hence $\sum_n \Delta_n(t)$ converges uniformly in $t$ almost surely.
Define
\[
B(t)=\sum_{n=0}^\infty \Delta_n(t).
\]
Then $B$ has almost surely continuous sample paths.

\Step{4: Gaussianity and covariance structure}

For each fixed $t$, $B_N(t)$ is a finite linear combination of independent Gaussian
random variables, hence Gaussian. Moreover, $B_N(t)\to B(t)$ in $L^2$, so $B(t)$ is
Gaussian.

For $s,t\in[0,1]$,
\[
\E[B(s)B(t)]
=
\sum_{n,k} f_{k,n}(s)f_{k,n}(t).
\]
Since
\[
f_{k,n}(t)=\int_0^t g_{k,n}(u)\,du
= \langle 1_{[0,t]}, g_{k,n}\rangle_{L^2},
\]
the bilinear Parseval identity yields
\[
\E[B(s)B(t)]
=
\langle 1_{[0,s]},1_{[0,t]}\rangle_{L^2}
=
\int_0^{\min(s,t)}1\,du
=
\min(s,t).
\]

\Step{5: Identification as Brownian motion}

Since $(B(s),B(t))$ is jointly Gaussian,
\[
B(t)-B(s)\sim\mathcal N(0,t-s).
\]
Moreover, for disjoint intervals $[s,t]$ and $[u,v]$,
\[
\mathrm{Cov}(B(t)-B(s),\,B(v)-B(u))=0,
\]
which implies independence of increments.
Finally, $B_0=0$ holds trivially from the definition.

Therefore $B$ is a standard Brownian motion on $[0,1]$.

\end{proof}

\begin{remark}
This construction provides both a theoretical realization of Brownian motion
and a practical approximation scheme: truncating the series yields continuous
random paths converging uniformly almost surely.
\end{remark}
